% !TEX root = ../main.tex

\begin{digest}
  Anterior cruciate ligament (ACL) tears and multi-ligament knee injuries (MLKI) are common and clinically important causes of post-traumatic osteoarthritis (PTOA). Although both conditions can lead to long-term degeneration of the knee joint, their injury patterns and early imaging manifestations are different. Isolated ACL tears are often associated with pivot-shift trauma and focal impaction between the femoral condyle and tibial plateau, whereas MLKI is usually caused by higher-energy trauma and involves more extensive disruption of the stabilizing structures of the knee. Magnetic resonance imaging (MRI) provides good soft-tissue contrast and is the main imaging modality for evaluating cartilage, bone, joint effusion, and bone marrow lesions (BML) after knee trauma. Quantitative three-dimensional analysis of cartilage morphology and BML extent has the potential to support early risk assessment, treatment planning, and understanding of PTOA risk.

  Accurate quantitative analysis of knee MRI remains difficult in real clinical scenarios. Pixel-wise or voxel-wise manual annotation of knee structures is expensive and time-consuming. The articular cartilage is thin, curved, and spatially complex, and BML often has diffuse and ambiguous boundaries. Manual segmentation of a single three-dimensional knee MRI scan can require considerable expert effort and may still suffer from inter-observer variability. MRI data acquired from different scanners and pulse sequences may also have large domain shifts. Public datasets such as Osteoarthritis Initiative-Zuse Institute Berlin (OAI-ZIB) provide high-resolution three-dimensional double echo steady-state (DESS) images with high-quality bone and cartilage labels, but many clinical trauma cohorts are acquired using routine fat-suppressed fast spin echo (FSE) sequences and do not contain expert segmentation labels. Models trained only on the labeled source domain may generalize poorly to the unlabeled clinical target domain. This thesis therefore focuses on unsupervised domain adaptation for bone and cartilage segmentation, validation on a small manually annotated target-domain test set, automated quantification of candidate BML high-signal volume, and quantitative comparison between isolated ACL tears and MLKI.

  For the segmentation component, this work constructs a cross-sequence framework based on unsupervised domain adaptation. The labeled source domain consists of the OAI-ZIB dataset, which contains \num{507} sagittal DESS knee MRI volumes and expert-corrected labels for femoral bone, femoral cartilage, tibial bone, and tibial cartilage. The target domain consists of clinical fat-suppressed FSE images collected from patients with isolated ACL tears and MLKI. To reduce scanner- and sequence-related variation, the clinical Digital Imaging and Communications in Medicine (DICOM) data were converted to Neuroimaging Informatics Technology Initiative (NIfTI) format, corrected using N4 bias field correction, resampled to a consistent physical spacing, cropped or padded to a common in-plane matrix size, and normalized using three-dimensional foreground-based Z-score standardization. A small independent FSE test set with manual annotations was used only for performance evaluation.

  The segmentation network uses an improved Domain Adversarial Neural Network (DANN). The network includes a ResNet34-based encoder, a U-Net-like decoder, and a domain discriminator connected through a gradient reversal layer. Batch normalization layers in the backbone were replaced with instance normalization to reduce the influence of image-level contrast and intensity style. An Atrous Spatial Pyramid Pooling (ASPP) module was inserted near the bottleneck to enlarge the receptive field and strengthen multi-scale contextual representation for thin cartilage structures. A dynamic freezing strategy based on the ratio between segmentation loss and adversarial loss was also designed to mitigate instability in adversarial training. When the discriminator became too strong, its parameters were temporarily frozen so that the feature extractor could continue learning more domain-confusing representations; when the discriminator became too weak, the feature extractor was temporarily frozen to restore the balance of the adversarial game. Because a complete ablation study was not performed, these architectural choices should be interpreted as design hypotheses supported by the model's target-domain performance rather than independently proven sources of improvement.

  The target-domain segmentation results show that the improved DANN achieved Dice scores of $0.9557 \pm 0.0243$ for femoral bone, $0.7748 \pm 0.0633$ for femoral cartilage, $0.9678 \pm 0.0231$ for tibial bone, and $0.7530 \pm 0.1505$ for tibial cartilage. The corresponding average surface distances were \qty{1.039}{\milli\metre}, \qty{0.961}{\milli\metre}, \qty{0.547}{\milli\metre}, and \qty{0.678}{\milli\metre}, respectively. These results indicate that the model obtained usable bone and cartilage segmentations on unlabeled clinical FSE data, especially for large bony structures. However, the experiment mainly evaluated a single improved DANN configuration and did not provide systematic quantitative comparisons against Source-only U-Net, standard DANN, DANN+IN, DANN+ASPP, or DANN with dynamic freezing. Cartilage segmentation remained more challenging because of its small thickness, complex curvature, and reduced contrast in clinical FSE images. The relatively large maximum surface distance in femoral cartilage also suggests that local boundary errors and isolated false positives still exist. Therefore, the segmentation results are suitable for exploratory quantitative analysis, but further validation, baseline comparison, and ablation studies are still needed before direct clinical deployment.

  For quantitative analysis, automatic tools were developed for cartilage morphology and joint spatial relationship assessment. Based on the predicted bone and cartilage masks, cartilage thickness was estimated using a three-dimensional Euclidean distance transform. The bone-cartilage interface and articular surface were extracted, and the local distance between these surfaces was used as a thickness estimate. A mesh-based image-derived near-contact area calculation was also implemented. Femoral and tibial cartilage masks were converted to triangular meshes using marching cubes. A kd-tree structure was then used to search for femoral cartilage surface elements within \qty{1}{\milli\metre} of the tibial cartilage surface. These elements were defined as near-contact regions under the MRI acquisition posture. Because routine MRI is acquired in a supine and non-weight-bearing position, this metric should be interpreted as an image-based spatial proximity measure rather than a direct measurement of true load-bearing cartilage contact.

  The resulting workflow was applied to a clinical cohort to compare isolated ACL tears and MLKI. The original clinical dataset included \num{500} isolated ACL cases and \num{165} MLKI cases. After image quality control and segmentation confidence screening, \num{430} isolated ACL cases and \num{149} MLKI cases, for a total of \num{579} patients, were included in the statistical analysis. Welch's two-sample t-test was first used as a reference analysis for mean differences because the extracted image-derived features may have unequal variances between groups. However, the image-based tibiofemoral near-contact area contained many zero values and showed a skewed distribution, especially in the MLKI group; therefore, its interpretation should also consider median [IQR], boxplot morphology, and robustness analyses rather than relying only on means. The MLKI group showed a smaller mean image-based tibiofemoral near-contact area than the isolated ACL group, with mean values of \qty{125.68}{\square\milli\metre} and \qty{244.54}{\square\milli\metre}, respectively ($t=7.00$, $p=1.67\times10^{-11}$). In age- and sex-adjusted regression, MLKI was still associated with a lower near-contact area ($\beta=\qty{-110.78}{\square\milli\metre}$, 95\%CI: \qtyrange{-148.01}{-73.55}{\square\milli\metre}, $p=8.53\times10^{-9}$). The MLKI group also showed a larger candidate BML high-signal volume than the isolated ACL group, with mean values of \qty{5217.77}{\cubic\milli\metre} and \qty{3907.14}{\cubic\milli\metre}, respectively ($t=-4.91$, $p=1.74\times10^{-6}$); this difference remained after age and sex adjustment ($\beta=\qty{1326.82}{\cubic\milli\metre}$, 95\%CI: \qtyrange{813.58}{1840.05}{\cubic\milli\metre}, $p=5.17\times10^{-7}$).

  These findings suggest that isolated ACL tears and MLKI differ in quantitative MRI-derived joint spatial relationships and bone marrow edema extent. The lower image-based near-contact area in MLKI may be related to ligamentous disruption, residual joint instability, joint effusion, or altered joint alignment under the non-weight-bearing MRI posture. The larger candidate BML high-signal volume in MLKI is consistent with the higher-energy and more complex trauma patterns often involved in multi-ligament injuries. However, these results should be interpreted as observational associations rather than direct causal proof of a specific injury process. Important confounders such as time from injury to MRI, effusion volume, injury severity, reduction status, scanning parameters, and longitudinal clinical outcomes were not fully modeled in this thesis.

  Overall, this thesis presents an unsupervised domain adaptation framework for knee MRI bone and cartilage segmentation, and develops automatic tools for three-dimensional cartilage thickness mapping, image-based near-contact area estimation, and candidate BML high-signal volume analysis. The workflow supports quantitative comparison between isolated ACL tears and MLKI without requiring full manual annotation of all clinical scans. The results provide preliminary quantitative evidence that these two injury patterns differ in joint spatial proximity and candidate BML burden. A multi-view high-resolution reconstruction direction based on rigid registration, B-spline refinement, and Deep Image Prior (DIP) was also explored during the study, but it was retained as a limitation and future-work item because systematic quantitative validation and downstream statistical use were not completed. Future work should expand the manually annotated target-domain test set, complete baseline and ablation experiments for the DANN variants, improve three-dimensional segmentation continuity, quantitatively validate DIP-based reconstruction, and integrate clinical follow-up data to better understand the progression from acute trauma to PTOA.
\end{digest}
